% Options for packages loaded elsewhere
\PassOptionsToPackage{unicode}{hyperref}
\PassOptionsToPackage{hyphens}{url}
\documentclass[
  ignorenonframetext,
]{beamer}
\newif\ifbibliography
\usepackage{pgfpages}
\setbeamertemplate{caption}[numbered]
\setbeamertemplate{caption label separator}{: }
\setbeamercolor{caption name}{fg=normal text.fg}
\beamertemplatenavigationsymbolsempty
% remove section numbering
\setbeamertemplate{part page}{
  \centering
  \begin{beamercolorbox}[sep=16pt,center]{part title}
    \usebeamerfont{part title}\insertpart\par
  \end{beamercolorbox}
}
\setbeamertemplate{section page}{
  \centering
  \begin{beamercolorbox}[sep=12pt,center]{section title}
    \usebeamerfont{section title}\insertsection\par
  \end{beamercolorbox}
}
\setbeamertemplate{subsection page}{
  \centering
  \begin{beamercolorbox}[sep=8pt,center]{subsection title}
    \usebeamerfont{subsection title}\insertsubsection\par
  \end{beamercolorbox}
}
% Prevent slide breaks in the middle of a paragraph
\widowpenalties 1 10000
\raggedbottom
\AtBeginPart{
  \frame{\partpage}
}
\AtBeginSection{
  \ifbibliography
  \else
    \frame{\sectionpage}
  \fi
}
\AtBeginSubsection{
  \frame{\subsectionpage}
}
\usepackage{iftex}
\ifPDFTeX
  \usepackage[T1]{fontenc}
  \usepackage[utf8]{inputenc}
  \usepackage{textcomp} % provide euro and other symbols
\else % if luatex or xetex
  \usepackage{unicode-math} % this also loads fontspec
  \defaultfontfeatures{Scale=MatchLowercase}
  \defaultfontfeatures[\rmfamily]{Ligatures=TeX,Scale=1}
\fi
\usepackage{lmodern}
\ifPDFTeX\else
  % xetex/luatex font selection
\fi
% Use upquote if available, for straight quotes in verbatim environments
\IfFileExists{upquote.sty}{\usepackage{upquote}}{}
\IfFileExists{microtype.sty}{% use microtype if available
  \usepackage[]{microtype}
  \UseMicrotypeSet[protrusion]{basicmath} % disable protrusion for tt fonts
}{}
\makeatletter
\@ifundefined{KOMAClassName}{% if non-KOMA class
  \IfFileExists{parskip.sty}{%
    \usepackage{parskip}
  }{% else
    \setlength{\parindent}{0pt}
    \setlength{\parskip}{6pt plus 2pt minus 1pt}}
}{% if KOMA class
  \KOMAoptions{parskip=half}}
\makeatother
\setlength{\emergencystretch}{3em} % prevent overfull lines
\providecommand{\tightlist}{%
  \setlength{\itemsep}{0pt}\setlength{\parskip}{0pt}}
\usepackage{bookmark}
\IfFileExists{xurl.sty}{\usepackage{xurl}}{} % add URL line breaks if available
\urlstyle{same}
\hypersetup{
  hidelinks,
  pdfcreator={LaTeX via pandoc}}

\author{\texorpdfstring{}{}}
\date{}

\begin{document}

\begin{frame}[fragile]{Appendix C --- Galois Connection Proof Sketch}
\protect\phantomsection\label{appendix-c-galois-connection-proof-sketch}
This appendix gives the formal statement and proof sketch of the
ε-approximate Galois connection between the category of Python program
graphs and the category of GNN generative models. The informal version
appears in Section 6.2 of the main text.

\begin{block}{C.1 Categories}
\protect\phantomsection\label{c.1-categories}
Let \textbf{Prog} be the category whose objects are typed Python program
graphs \texttt{G\ =\ (V,\ E,\ λ\_V,\ λ\_E,\ τ)} in the sense of Section
2.2 (14 node kinds, 11 edge kinds) and whose morphisms are graph
homomorphisms that preserve node and edge labels. Let \textbf{GNN} be
the category whose objects are GNN v1.1 bundles (the Markdown sections
\texttt{StateSpaceBlock}, \texttt{Connections},
\texttt{InitialParameterization}, \texttt{ActInfOntologyAnnotation},
plus the A/B/C/D matrices) and whose morphisms are role-preserving
bundle embeddings.

Both categories are posets under the pointwise subset order:
\texttt{G\ ≤\ G\textquotesingle{}} in \textbf{Prog} iff
\texttt{V\ ⊆\ V\textquotesingle{}}, \texttt{E\ ⊆\ E\textquotesingle{}},
and the labelings agree on the common subset;
\texttt{M\ ≤\ M\textquotesingle{}} in \textbf{GNN} iff each bundle
section of \texttt{M} is included in the corresponding section of
\texttt{M\textquotesingle{}}.
\end{block}

\begin{block}{C.2 Forward and reverse functors}
\protect\phantomsection\label{c.2-forward-and-reverse-functors}
Define two order-preserving maps:

\begin{quote}
\textbf{F : Prog → GNN} --- the forward pipeline. \texttt{F(G)} is the
GNN bundle emitted by \texttt{cogant\ translate\ G}: it runs ingest →
static → normalize → graph → translate → statespace → process → export →
validate and returns the \texttt{gnn\_package/model.gnn.md} bundle
together with the derived A/B/C/D matrices.

\textbf{R : GNN → Prog} --- the reverse pipeline. \texttt{R(M)} is the
typed program graph extracted by running \texttt{cogant\ reverse\ M}
(which internally invokes
\texttt{parse\_gnn\ →\ plan\_package\ →\ synthesize\_package}) and then
re-parsing the synthesized Python package through the static + graph
stages of the forward pipeline.
\end{quote}

Both \texttt{F} and \texttt{R} are monotone because each underlying
stage is monotone: adding a node or edge to the input graph can only add
mappings to \texttt{semantic\_mappings.json}, which can only add
declarations to the GNN bundle, which can only add planned nodes to
\texttt{plan\_package}, which can only add synthesized code artefacts.
\end{block}

\begin{block}{C.3 Role multiset functor}
\protect\phantomsection\label{c.3-role-multiset-functor}
Define the role-multiset functor \textbf{ρ : Prog → Mset(Roles)} that
sends a program graph \texttt{G} to the multiset of Active Inference
roles assigned to its nodes by the translate engine, where
\texttt{Roles\ =\ \{HIDDEN\_STATE,\ OBSERVATION,\ ACTION,\ POLICY,\ CONSTRAINT,\ CONTEXT\}}.
Extend \texttt{ρ} to \textbf{GNN → Mset(Roles)} by counting the
declarations in each role-tagged section of a GNN bundle
(\texttt{StateSpaceBlock} → HIDDEN\_STATE, observation modalities →
OBSERVATION, control states → ACTION, etc.). Both extensions agree on
the image of \texttt{F}: \texttt{ρ(F(G))\ =\ ρ\_GNN(F(G))} for every
\texttt{G\ ∈\ Prog}, because the forward pipeline emits one section
entry per mapping in the translate output.
\end{block}

\begin{block}{C.4 Adjunction (approximate)}
\protect\phantomsection\label{c.4-adjunction-approximate}
\textbf{Proposition C.1 (ε-approximate Galois connection).} The
forward/reverse pair \texttt{(F,\ R)} satisfies the approximate Galois
condition

\begin{quote}
**F(G) ≤\_GNN M ⟺\_ε G ≤\_Prog R(M)**
\end{quote}

where \texttt{⟺\_ε} means ``the two inequalities agree on at least the
ε-fraction of the role multiset'', i.e.~for every \texttt{G\ ∈\ Prog}
and \texttt{M\ ∈\ GNN}:

\begin{quote}
\texttt{multiset\_sim(ρ(G),\ ρ(R(F(G))))\ ≥\ 1\ −\ ε\_worst}
\end{quote}

where \texttt{ε\_worst} depends only on the rule table and the
synthesizer.

\textbf{Proof sketch.} The forward pipeline is the composition of a
finite sequence of monotone rule applications (the 19 translation rules
in the translate engine plus the A/B/C/D derivation in
\texttt{statespace}), each of which emits exactly one mapping per
triggering graph pattern. The reverse pipeline is the composition of
\texttt{parse\_gnn} (which is a right inverse of the GNN emitter by
construction --- the emitter's output is parseable by its own parser)
and \texttt{synthesize\_package} (which emits one Python function per
planned node). Composing the two:

\begin{enumerate}
\tightlist
\item
  Start with \texttt{G\ ∈\ Prog}.
\item
  \texttt{F(G)} emits one GNN declaration per mapping in
  \texttt{translate(G)}; the number of declarations of role \texttt{r}
  equals \texttt{count\_ρ(G,\ r)}.
\item
  \texttt{parse\_gnn(F(G))} recovers the full declaration list
  bijectively.
\item
  \texttt{synthesize\_package(plan)} emits one Python artefact per
  \texttt{NodePlan}; by the wave‑14 CONSTRAINT fix, the mapping from
  \texttt{NodePlan} to emitted artefact is injective on role
  multiplicity.
\item
  Re‑running \texttt{F} on the synthesized package recovers the same
  role multiset up to the \textbf{synthesizer gap}: extra
  OBSERVATION/CONSTRAINT nodes produced by scaffolding, which inflate
  \texttt{count\_ρ(R(F(G)),\ r)} for those roles but preserve the origin
  roles exactly.
\end{enumerate}

The multiset similarity \texttt{min(a,b)\ /\ max(a,b)} averaged over
roles is therefore bounded below by
\texttt{(count\_origin)\ /\ (count\_origin\ +\ scaffold\_r)} for each
role \texttt{r}, where \texttt{scaffold\_r} is the fixed contribution of
the reverse synthesizer's scaffolding (4 CONSTRAINT, 7 OBSERVATION, 5
ACTION on the minimum-case synthesis). The worst-case ε is achieved on
targets where the origin role counts are smaller than the scaffold; on
zoo fixtures (small origin, scaffold dominates) the ratio saturates
because both sides collapse to the scaffold, and on real-world libraries
(large origin, scaffold negligible) the ratio approaches 1.0 once the
CONSTRAINT fix is applied. In both regimes the Galois condition holds up
to a bounded ε that depends only on the rule table and the fixed
synthesizer scaffolding. ∎
\end{block}

\begin{block}{C.5 ε-isomorphism theorem}
\protect\phantomsection\label{c.5-ux3b5-isomorphism-theorem}
\textbf{Theorem C.2 (ε-Isomorphism).} For any \texttt{P\ ∈\ Prog}, the
roundtrip \texttt{P\ →\ F(P)\ →\ R(F(P))} preserves the role
distribution up to

\begin{quote}
\textbf{ε(P, R(F(P))) = JS(ρ\_norm(P) ∥ ρ\_norm(R(F(P))))}
\end{quote}

where \texttt{ρ\_norm} is the role multiset normalized to a probability
distribution over \texttt{Roles}, and \texttt{JS} is the Jensen--Shannon
distance. When the multiset-similarity implementation in
\texttt{compute\_isomorphism\_report.role\_match\_score} is substituted
for \texttt{JS}, the theorem holds with the multiset-similarity metric
in place of JS‑distance and yields the values reported in Appendix A.

\textbf{Proof sketch.} The translate engine emits one
\texttt{SemanticMapping} per triggered rule, and each mapping carries
exactly one role label. The forward GNN bundle's
\texttt{StateSpaceBlock}, observation modalities, control states, and
constraint annotations are in one-to-one correspondence with those
mappings, so \texttt{ρ\_norm(F(P))\ =\ ρ\_norm(P)} (the forward map is
role-preserving up to normalization). The reverse map introduces
scaffolding nodes that inflate the role counts additively:
\texttt{count(R(F(P)),\ r)\ =\ count(P,\ r)\ +\ scaffold\_r} for each
role. The Jensen--Shannon distance between \texttt{P} and
\texttt{R(F(P))} is therefore bounded by the JS distance between two
distributions that differ only by a fixed additive shift, which in turn
is bounded by a function of
\texttt{sum\_r\ scaffold\_r\ /\ sum\_r\ count(P,\ r)}. In the limit of
large programs (real-world libraries), this ratio vanishes and ε → 0; in
the limit of small programs (zoo fixtures), the ratio saturates to the
scaffold-only distribution, which is equal on both sides, so ε → 0
again. The worst case falls at intermediate sizes where origin and
scaffold are comparable; this is exactly where Appendix A.1 shows
overall ε ≈ 0.85--0.95. ∎
\end{block}

\begin{block}{C.6 ISOMORPHIC threshold corresponds to majority role
preservation}
\protect\phantomsection\label{c.6-isomorphic-threshold-corresponds-to-majority-role-preservation}
\textbf{Proposition C.3.} The threshold \texttt{ε\ ≥\ 0.5} used
throughout the paper to classify a target as ISOMORPHIC is equivalent to
``a majority of the origin role multiset is preserved in the
roundtrip''.

\textbf{Proof.} The multiset similarity per role is
\texttt{min(a,b)\ /\ max(a,b)}. Averaging over the \texttt{k} roles
present on either side and requiring the mean ≥ 0.5 is equivalent to
requiring that at least \texttt{k/2} of the per-role ratios are ≥ 0.5
(or a weighted combination of more and fewer). For a single role,
\texttt{min(a,b)\ /\ max(a,b)\ ≥\ 0.5} iff
\texttt{max(a,b)\ ≤\ 2·min(a,b)} iff each count is at most twice the
other. When the reverse synthesizer only adds scaffolding, this is
equivalent to requiring \texttt{count\_origin\ ≥\ scaffold\ /\ 2},
i.e.~that the origin population is at least half of the synth
population. Summing over roles, the ISOMORPHIC threshold corresponds to
``the majority of the origin role multiset survives to the roundtrip
without being drowned out by scaffolding''. The CONSTRAINT fix (§A.2) is
exactly the transformation that makes this true for constraint-heavy
real-world libraries: it raises the CONSTRAINT component of
\texttt{count\_synth} from 3 (scaffolding) to \texttt{count\_origin}
(proportional), so \texttt{min\ =\ count\_origin} and the per-role ratio
jumps to 1.0. ∎
\end{block}
\end{frame}

\begin{frame}
\end{frame}

\end{document}
